% 贯穿图的子图与诱导子图：S={1,2,3}。虚线灰色表示被删除的顶点或边。
\begin{tikzpicture}[>=Stealth,
  v/.style={circle,draw=FigureStroke,fill=FigureBlueFill,line width=.8pt,minimum size=6.5mm,inner sep=0pt},
  vd/.style={circle,draw=FigureMuted,dashed,fill=white,text=FigureMuted,line width=.6pt,minimum size=6.5mm,inner sep=0pt},
  e/.style={draw=FigureStroke,line width=.8pt},
  ed/.style={draw=FigureMuted,dashed,line width=.6pt}]
\newcommand{\pos}{\coordinate (p1) at (-1.0,0); \coordinate (p2) at (0,.8);
  \coordinate (p3) at (0,-.8); \coordinate (p4) at (1.0,0);}
\begin{scope}
  \pos
  \node[v] (1) at (p1) {$1$}; \node[v] (2) at (p2) {$2$}; \node[v] (3) at (p3) {$3$}; \node[v] (4) at (p4) {$4$};
  \draw[e] (1)--(2) (1)--(3) (2)--(3) (2)--(4) (3)--(4);
  \node at (0,-1.5) {(a) 原图$\mathcal G$};
\end{scope}
\begin{scope}[xshift=2.8cm]
  \pos
  \node[v] (1) at (p1) {$1$}; \node[v] (2) at (p2) {$2$}; \node[v] (3) at (p3) {$3$}; \node[vd] (4) at (p4) {$4$};
  \draw[e] (1)--(2) (1)--(3) (2)--(3); \draw[ed] (2)--(4) (3)--(4);
  \node[align=center] at (0,-1.65) {(b) 删除顶点$4$：\\$\mathcal G-\{4\}=\mathcal G[\{1,2,3\}]$};
\end{scope}
\begin{scope}[xshift=5.6cm]
  \pos
  \node[v] (1) at (p1) {$1$}; \node[v] (2) at (p2) {$2$}; \node[v] (3) at (p3) {$3$}; \node[vd] (4) at (p4) {$4$};
  \draw[e] (1)--(2) (2)--(3); \draw[ed] (1)--(3) (2)--(4) (3)--(4);
  \node[align=center] at (0,-1.65) {(c) 同样三个顶点，\\缺少$\{1,3\}$：非诱导};
\end{scope}
\begin{scope}[xshift=8.4cm]
  \pos
  \node[v] (1) at (p1) {$1$}; \node[v] (2) at (p2) {$2$}; \node[v] (3) at (p3) {$3$}; \node[v] (4) at (p4) {$4$};
  \draw[e] (1)--(2) (1)--(3) (2)--(4) (3)--(4); \draw[ed] (2)--(3);
  \node[align=center] at (0,-1.65) {(d) 只删除边$\{2,3\}$，\\保留全部顶点};
\end{scope}
\end{tikzpicture}
